
\textbf{Immediate Extensions}: Real PWC data validation is the highest priority---re-running all experiments with actual Papers With Code leaderboard data when API access is restored will unlock real-world applicability claims and determine whether per-benchmark calibration is genuinely required or a synthetic artifact. Dual-metric integration should implement a combined detector requiring both score convergence and velocity decay, comparing precision/recall against single-metric baselines on an expanded benchmark set (n = 10+). Real expert survey collection should execute with IRB approval, targeting n = 100--150 responses via NeurIPS/ICML mailing lists to validate expert consensus with actual ML researcher opinions.

\textbf{Sample Size Expansion}: Expanding temporal precedence analysis from n = 3 to n = 15+ benchmark-shift pairs is necessary for statistical significance (current p = 0.125 $\rightarrow$ target p < 0.05). Additional pairs include CIFAR-$10\rightarrow ResNet, SuperGLUE\rightarrow T5, WikiText\rightarrow GPT-2$, MS $COCO\rightarrow Mask$ R-CNN.

\textbf{Citation Mechanism Refinement}: Testing aligned 9-month windows (both detector and ground truth) and threshold variants ($3\sigma, 4\sigma$ spike detection) may improve precision from 0.50 toward 0.80 target.

\textbf{Longer-Term Vision}: Developing community coordination mechanisms for systematic benchmark lifecycle management and piloting conference policy integration (saturation dashboard with NeurIPS Benchmark Track or ICLR) to test voluntary adoption hypothesis will track adoption rate (percentage of papers referencing saturation status) and author survey feedback.

The shortest path from proof-of-concept to production deployment is real PWC data validation, which determines whether mechanisms validated on synthetic data generalize to actual leaderboard submissions. This single validation unlocks applicability claims and enables infrastructure piloting with conference organizers.

Benchmark saturation is not an inevitable background process to be endured, but a detectable phenomenon with measurable signals and actionable lead times. Systematic saturation detection transforms benchmark lifecycle from unmanaged monument to governed consumable, enabling the ML research community to invest resources where genuine progress remains possible.


Bansal, G., Nushi, B., Kamar, E., Lasecki, W. S., Weld, D. S., \& Horvitz, E. (2021). Beyond accuracy: The role of mental models in human-AI team performance. \textit{Proceedings of the AAAI Conference on Human Computation and Crowdsourcing}, 9(1), 2-11.

Brown, T. B., Mann, B., Ryder, N., Subbiah, M., Kaplan, J., Dhariwal, P., ... \& Amodei, D. (2020). Language models are few-shot learners. \textit{Advances in Neural Information Processing Systems}, 33, 1877-1901.

Gebru, T., Morgenstern, J., Vecchione, B., Vaughan, J. W., Wallach, H., Daumé III, H., \& Crawford, K. (2018). Datasheets for datasets. \textit{arXiv preprint arXiv:1803.09010}.

Hendrycks, D., \& Dietterich, T. (2019). Benchmarking neural network robustness to common corruptions and perturbations. \textit{International Conference on Learning Representations}.

Nie, Y., Williams, A., Dinan, E., Bansal, M., Weston, J., \& Kiela, D. (2020). Adversarial NLI: A new benchmark for natural language understanding. \textit{Proceedings of the 58th Annual Meeting of the Association for Computational Linguistics}, 4885-4901.

Pineau, J., Vincent-Lamarre, P., Sinha, K., Larivière, V., Beygelzimer, A., d'Alché-Buc, F., ... \& Larochelle, H. (2021). Improving reproducibility in machine learning research (a report from the NeurIPS 2019 reproducibility program). \textit{Journal of Machine Learning Research}, 22(164), 1-20.
